\documentclass[11pt,a4paper]{article}
\usepackage[T1]{fontenc}
\usepackage{lmodern}
\usepackage[margin=30mm,headheight=14pt]{geometry}
\usepackage{amsmath,amssymb,microtype}
\usepackage{enumitem,fancyhdr,hyperref}
\hypersetup{pdftitle={Scientific Computing 2026: Seminar 1}}
\setlength{\parindent}{0pt}
\setlength{\parskip}{5pt}
\setlist[enumerate]{label=(\alph*),leftmargin=1.8em,itemsep=4pt,topsep=5pt}
\pagestyle{fancy}
\fancyhf{}
\fancyhead[L]{\small Seminar 1. Probability Models, Conditioning, and Random Variables}
\fancyhead[R]{\small\thepage}
\renewcommand{\headrulewidth}{0pt}
\newcommand{\Prob}{\mathbb P}
\newcommand{\problem}[2]{\par\addvspace{6pt}\textbf{Problem #1} (#2). }
\begin{document}
\hypersetup{pageanchor=false}
\begin{titlepage}
\centering
\vspace*{12mm}
{\Huge Scientific Computing\par}
\vspace{5mm}
{\large Mathematical Foundations for Machine Learning\par}
\vspace{24mm}
{\LARGE Seminar 1: Problem Set\par}
\vspace{4mm}
{\large Probability Models, Conditioning, and Random Variables\par}
\vspace{12mm}
{\large Block I, following Lecture 1\par}
\vfill
{\large 2026\par}
\end{titlepage}
\setcounter{page}{1}
\hypersetup{pageanchor=true}

\problem{1}{The Monty Hall protocol}
Three doors. Behind one of them is an internship in a decent laboratory;
behind the other two, goats. The prize door is uniform over the three doors,
independently of your initial choice. You point at a door. The host, who
knows what is behind each door, opens an unchosen door hiding a goat and
offers you the chance to switch. If both unchosen doors hide goats, he
chooses between them with equal probabilities. Write \(c\) for the prize
door, \(g\) for your initial guess, and \(h\) for the door the host opened.
\begin{enumerate}
\item Construct a probabilistic model of the experiment.
\item Assume you guessed door 1. Compute \(\Prob(c=1\mid h=3)\) and
\(\Prob(c=2\mid h=3)\) directly from the definition of conditional
probability. Should you switch?
\item Identify the exact place in your computation where the assumption
that the host knows the answer is used.
\end{enumerate}
\textbf{Remark.} This is the question put to a student in the film
\emph{21} (2008); answering it correctly earns him a seat on a card-counting
team. Evidence that a conditional probability exercise occasionally changes a life.

\problem{2}{Monty Fall}
The same three doors, but this time the host does not know what is behind
them. He trips, and in falling pushes open one of the two doors you did
not choose, each with probability \(1/2\), independently of the prize.
It happens to reveal a goat. He offers you the chance to switch.
\begin{enumerate}
\item Compute the probability that the prize is behind your original door,
given that a goat was revealed. Should you switch?
\item You see exactly what you saw in Problem 1: your door closed, one
other door open, a goat behind it. Explain why the answer is nevertheless different.
\item Which single quantity in the model of Problem 1 has changed?
\end{enumerate}

\problem{3}{A detector with 95\% sensitivity}
Consider a university that runs submitted essays through an AI-text
detector. Assume the detector has the following class-conditional probabilities:
\[
\Prob(\text{flag}\mid\text{AI})=0.95,\qquad
\Prob(\text{flag}\mid\text{human})=0.05.
\]
Assume one essay in twenty is AI-generated, so \(\Prob(\text{AI})=0.05\).
Your essay has been flagged.
\begin{enumerate}
\item Compute \(\Prob(\text{AI}\mid\text{flag})\).
\item How large would the base rate \(\Prob(\text{AI})\) have to be for
a single flag to give a posterior probability above \(0.9\)?
\end{enumerate}

\problem{4}{Two detectors, and the odds form of Bayes' rule}
The university adds a second detector with the same two characteristics,
and both detectors flag your essay. Write \(H\) for the hypothesis
``AI-generated'', \(H^c\) for its complement, and \(E_1,\ldots,E_n\) for
the events that detectors \(1,\ldots,n\) raise a flag. Assume that the
flags are mutually conditionally independent given each class.
\begin{enumerate}
\item Compute \(\Prob(H\mid E_1\cap E_2)\) directly from Bayes' theorem.
\item Define the posterior odds and prior odds
\[
\mathcal O(H\mid E)=\frac{\Prob(H\mid E)}{\Prob(H^c\mid E)},
\qquad \mathcal O(H)=\frac{\Prob(H)}{\Prob(H^c)},
\]
and the likelihood ratio of a single flag,
\(\mathrm{LR}=\Prob(E_i\mid H)/\Prob(E_i\mid H^c)\).
Write Bayes' theorem for \(H\) and \(H^c\), take the ratio, and cancel
the common normalizing denominator. Using conditional independence in
both classes, deduce that for \(n\) detectors with identical characteristics
\[
\mathcal O(H\mid E_1\cap\cdots\cap E_n)
=\mathrm{LR}^{n}\,\mathcal O(H).
\]
Convert back with \(P=\mathcal O/(1+\mathcal O)\), check that \(n=2\)
reproduces part (a), and find the smallest \(n\) for which the posterior
probability exceeds \(0.99\).
\end{enumerate}
\textbf{Remark (Dependent flags).} Without conditional independence,
replace \(\mathrm{LR}^{n}\) by the actual joint likelihood ratio
\[
\frac{\Prob(E_1\cap\cdots\cap E_n\mid H)}
     {\Prob(E_1\cap\cdots\cap E_n\mid H^c)}.
\]
Dependence alone does not tell us whether this ratio is above or below
the product of the individual ratios: the joint probabilities in both
classes determine it. If the second detector always copies the first,
then \(E_2=E_1\), and two flags give exactly the same posterior as one.
Multiplying the two likelihood ratios would count the same evidence twice.

\problem{5}{The best of several runs}
A random search launches two independent training runs. Their qualities
\(R_1,R_2\) are uniform on \(\{1,2,\ldots,6\}\) on a fixed scale.
You keep the better of the two: \(X=\max(R_1,R_2)\).
\begin{enumerate}
\item Find the CDF \(F_X(k)\) for \(k=1,\ldots,6\) before attempting
the PMF, and explain why the CDF is easier to calculate here.
\item Recover the PMF using \(p_X(k)=F_X(k)-F_X(k^-)\).
\item Compute \(\Prob(X\ge5)\).
\item Generalize to \(n\) independent runs. How large must \(n\) be for
\(\Prob(X=6)>0.9\)?
\end{enumerate}

\problem{6}{Densities and inverse-CDF sampling}
\begin{enumerate}
\item Let \(f_Y(y)=10\) for \(y\in[0,0.1]\), and zero otherwise.
Verify that \(f_Y\) is a density and compute \(\Prob(Y=0.05)\).
\item Let \(f_X(x)=2x\) on \([0,1]\), and zero elsewhere. Find \(F_X\)
and its inverse \(F_X^{-1}\) on \([0,1]\).
\item Let \(U\sim\operatorname{Uniform}(0,1)\). Show that
\(X=F_X^{-1}(U)\) has CDF \(F_X\). Begin with \(\Prob(X\le x)\).
\item Write down the same construction for the exponential distribution
with \(F(x)=1-e^{-\lambda x}\), \(x\ge0\).
\end{enumerate}
\end{document}
